

\section{Affine Algebraic Geometry}
The object of this section is to give a basic \textbf{geometry-algebra correspondence} between affine algebraic varieties and a polynomial ring over a field, which can be concluded by the following graph:
% https://q.uiver.app/#q=WzAsNCxbMCwwLCJrXm4iXSxbMywwLCJcXG1hdGhjYWx7T30oa15uKSJdLFszLDMsImtbWF8xLC4uLixYX25dIl0sWzAsMywiXFxtYXRoYmJ7QX1ebl9rIl0sWzAsMSwiXFxtYXRoY2Fse099Il0sWzIsMSwiXFxtYXRocm17Q2hhcn0oayk9MCIsMix7InN0eWxlIjp7InRhaWwiOnsibmFtZSI6ImFycm93aGVhZCJ9LCJib2R5Ijp7Im5hbWUiOiJkYXNoZWQifX19XSxbMiwzLCJcXG1hdGhybXtTcGVjfSJdLFswLDMsIlxcdGV4dHtOdWxsc3RlbGxlbnNhdHp9IiwyLHsic3R5bGUiOnsidGFpbCI6eyJuYW1lIjoiYXJyb3doZWFkIn0sImJvZHkiOnsibmFtZSI6ImRhc2hlZCJ9fX1dXQ==
\[\begin{tikzcd}
	{k^n} &&& {\mathcal{O}(k^n)} \\
	\\
	\\
	{\mathbb{A}^n_k} &&& {k[X_1,...,X_n]}
	\arrow["{\mathcal{O}}", from=1-1, to=1-4]
	\arrow["{\text{Nullstellensatz}}"', dashed, tail reversed, from=1-1, to=4-1]
	\arrow["{\text{infinite field}}"', dashed, tail reversed, from=4-4, to=1-4]
	\arrow["{\mathrm{Spec}}", from=4-4, to=4-1]
\end{tikzcd}\]





















\newpage


\subsection{Hilbert's zero theorem (Nullestellensatz)}

A classic question in old algebraic geometry is to study the solution of the polynomial equations over some field \(k\), for example the famous Fermat's last theorem is to study the integer solutions of the equation
\[x^n + y^n = z^n\]
if we see it as a polynoimial equation over the field \(\qq\). generally, if we have a set of polynoimial equations
\begin{equation}\label{eq:poly-equations}
    \begin{cases}
    f_1(x_1,\dots,x_n) = 0\\
    f_2(x_1,\dots,x_n) = 0\\
    \vdots \\
    f_m(x_1,\dots,x_n) = 0
    \end{cases}
\end{equation}
then its solution set is a subset of \(k^n\) which is called algebraic subset.Firstly we should clarify the isomorphism of polynomial functions and the formal polynomial in certain condition, which allows us to freely use to two objects without confusion.
\begin{definition}
    A formal polynomial is a element of the polynomial ring \(k[X_1,\dots,X_n]\) over a field \(k\), while a polynomial function is a function with a formal polynomial expression, i.e. a map
    \[f: k^n \to k, \quad (x_1,...,x_n) \mapsto \sum_{i=0}^{n}a_i x_i\]
    and we denote the set of all polynomial functions on \(k^n\) by \(\mathcal{O}(k^n)\).
\end{definition}

It is easy to verify that \(\mathcal{O}(k^n)\) is a \(k\)-algebra type fini, hence by the universal property of polynomial, there exists a unique surjective \(k\)-linear map
\[\mathrm{ev}:k[X_1,\dots,X_n] \to \mathcal{O}(k^n), \quad X_i \mapsto \pi_i: (x_1,...,x_n) \mapsto x_i\]
In particular, when \(k\) is an \textbf{infinite field}\textbf{ (verify and find counterexample)}, the map \(\mathrm{ev}\) is also injective, hence it is an isomorphism of \(k\)-algebras, which allows us to abuse the notation of polynomial functions and polynomials.
\begin{colordefinition}
    Let \(k\) be a field, and \(I\) is a ideal (subset) of the polynomial ring \(\oo(k^n)\), then we define the \textbf{algebraic subset} or \textbf{variety} of \(k^n\) defined by \(I\) to be
    \[\mathcal{V}(I) := \{(a_1,\dots,a_n) \in k^n \mid f(a_1,\dots,a_n) = 0, \forall f \in I\}\]
    it is also called the \textbf{variety} defined by \(I\).
\end{colordefinition}

\begin{remark} \label{varietymakestop}
    we should remark that the following facts\textbf{(verify!)}: \\
\[\mathcal{V}(S) = \mathcal{V}((S)), \quad \text{(S) is the ideal generated by S}\]
so the common zeros of \(f_1,\dots,f_m\) is the samething as the variety \(\mathcal{V}(I)\) defined by the ideal \(I = (f_1,\dots,f_m)\), so there is no problem to only consider the case of ideals. \\
\end{remark}

\begin{remark}
    The subset of \(k^n\) can not be a variety, for example \(n = 1\) and subset \(\zz \subset \rr\), If \(\zz = \mathcal{V}(I)\) for some ideal \(I \subset \rr[X]\), then there exists a polynomial \(f\) having infinitely many zeros, so it is impossible.
\end{remark}

From now on, if no clarification is given, the condition \(\mathcal{I} \subset \oo(k^n)\) will be assumed as an ideal. Moreover, some algebraic facts should be known:

\begin{proposition} \label{algberaic-subset}
    Let \(k\) be a field, then 
    \begin{enumerate}[(a)]
        \item  \(\mathcal{V}(1) = \emptyset\) and \(\mathcal{V}(0) = k^n\) as two trivial cases.
        \item \textcolor{blue}{(order-reversing)} If \(I \subset J\), then \(\mathcal{V}(I) \supset \mathcal{V}(J)\).
        \item \(\mathcal{V}(I)\cup \mathcal{V}(J) = \mathcal{V}(IJ) = \mathcal{V}(I \cap J)\) for any two ideals.
        \item \(\cap_{k}\mathcal{V}(I_k) = \mathcal{V}(\sum_k I_k)\) for any index.
        \item \(\mathcal{V}(I) = \mathcal{V}(\sqrt{I})\).
    \end{enumerate}
\end{proposition}

\begin{proof}
    
\end{proof}

Conversly an inverse problem can be given as following: for any subset \(V \subset k^n\), what is the largest set of polynomials that vanish on \(V\)? Which leads to the definition of \textbf{vanishing ideal}:
\[\mathcal{I}(V) = \{f \in \oo(k^n) \mid f|_V = 0 \}\]
equivalently, when \(k\) is an infinite field, we have
\[\mathcal{I}(V) = \{f \in k[X_1,\cdots,X_n] \mid f(a) = 0 ,  \forall a \in V\}\]

\begin{remark}
    \(I(V)\) is \textcolor{blue}{a radical ideal} since if \(f^n \in I(V)\), then \(f^n(a) = 0\) for any \(a \in V\), hence \(f(a) = 0\) for any \(a \in V\), which shows that \(f \in I(V)\), and review that 
    \[\text{maximal ideal} \subset \text{prime ideal} \subset \text{radical ideal}\]
    Moreover, \(\mathcal{I}(V)\) is exactly \textcolor{blue}{the largest ideal} of polynomials vanishing on \(V\): for any ideal \(J\) with \(\mathcal{V}(J) = V\), we have \(J \subset \mathcal{I}(V)\).
\end{remark}

\begin{proposition} \label{vanish-ideal}
    Let \(k\) be a field, then 
    \begin{enumerate}[(a)]
        \item \(\mathcal{I}(k^n) = \{0\}\) and \(\mathcal{I}(\emptyset) = \oo(k^n)\) as two trivial cases.
        \item \textcolor{blue}{(order-reversing)} If \(V \subset W\), then \(\mathcal{I}(V) \supset \mathcal{I}(W)\).
        \item \(\mathcal{I}(V)\cap \mathcal{I}(W) = \mathcal{I}(V\cup W)\) for any two subsets.
        \item \(\mathcal{I}(\cap_k V_k) = \sqrt{\sum_k \mathcal{I}(V_k)}\) for any index.
        \item \(\mathcal{V}(\mathcal{I}(S)) = S\) for any algebraic subset \(S \subset k^n\).
        \item \(J \subset \sqrt{J} \subset \mathcal{I}( \mathcal{V}(J)) \) for any ideal \(J \subset \oo(k^n)\).
    \end{enumerate}
\end{proposition}

\begin{proof}
        (4) \(V= \mathcal{V}(\mathcal{I}(V)) \) for any subset \(V \subset k^n\):\\

        It is clear that \(V \subset \mathcal{V}(\mathcal{I}(V))\), for avoid confusion of notation in traditional proof, we define a set \(\mathcal{J}=\{\mathcal{I} \subset K[X_1,\dots,X_n] \\\text{ ideal } \mid f|_V = 0, \forall f \in \mathcal{I}\}\), then the set is a partial order set under inclusion, it is easy to verify that \(\mathcal{I}(V)\) is a maximal element in it, i.e. each chain in \(\mathcal{J}\) will stop at \(\mathcal{V}(\mathcal{I}(V))\):
        \[\mathcal{I}_1 \subset \mathcal{I}_2 \subset \dots \subset \mathcal{I}(V)\]
        then immediately apply \(\mathcal{V}(-)\), we can get a decrasing chain of algebraic subsets:
        \[V \subset \mathcal{V}(\mathcal{I}(V)) \subset \dots \subset \mathcal{V}(\mathcal{I}_2) \subset \mathcal{V}(\mathcal{I}_1)\]
        \(\mathcal{V}(\mathcal{I}(V))\) must be the minimal element in the chain, hence \(\mathcal{V}(\mathcal{I}(V)) = V\).\\
\end{proof}

\begin{remark}
    The pair of maps can be drawn as following:
    % https://q.uiver.app/#q=WzAsNCxbMCwwLCJcXHRleHR7aWRlYWx9Il0sWzEsMSwiXFx0ZXh0e3N1YnNldCBvZiB9IGtebiJdLFsxLDAsIlxcdGV4dHtzdWJzZXQgb2YgfSBrXm4iXSxbMCwxLCJcXHRleHR7cmFkaWNhbCBpZGVhbH0iXSxbMCwyLCJcXG1hdGhjYWx7Vn0iXSxbMiwxLCIiLDAseyJzdHlsZSI6eyJ0YWlsIjp7Im5hbWUiOiJob29rIiwic2lkZSI6ImJvdHRvbSJ9fX1dLFszLDAsIiIsMCx7InN0eWxlIjp7InRhaWwiOnsibmFtZSI6Imhvb2siLCJzaWRlIjoiYm90dG9tIn19fV0sWzEsMywiXFxtYXRoY2Fse0l9Il1d
\[\begin{tikzcd}
	{\text{ideal}} & {\text{alg subset of } k^n} \\
	{\text{radical ideal}} & {\text{subset of } k^n}
	\arrow["{\mathcal{V}}", from=1-1, to=1-2]
	\arrow[hook', from=1-2, to=2-2]
	\arrow[hook', from=2-1, to=1-1]
	\arrow["{\mathcal{I}}", from=2-2, to=2-1]
\end{tikzcd}\]
(e) shows that \(\mathcal{I}\) is injective on the domain of algebraic subsets, while (f) shows that \(\mathcal{V}\) behaves not good.

Example 1: Let \(J = (X^2)\) then we can get \(\sqrt{J} = \mathcal{I}(\mathcal{V}(J))= (X)\), so we may guess that \(\mathcal{V}\) can be injective on the domain of radical ideals.

Example 2: Let \(k = \rr\) and \(J = (X^2+1)\) be a radical ideal, then we can calculate that \(\mathcal{V}(J) = \emptyset = \mathcal{V}(1)\), it shows that \(\mathcal{V}\) can not distinguish the radical ideal in \(\rr[X]\), so some restriction anout field is needed to make it work.


\end{remark}

Here we need a strong condition of field such that we can get a correpsondence betwween algebraic subsets and vanishing ideals, which motivates the Hilbert's zero theorem (Nullstellensatz), we review the original statement here:
\begin{theorem}[1893] \label{orginal-Nullstellensatz}
    Let \(f_1,...,f_m\) be the polynomials of \(\mathcal{O}(\cc^n)\), and \(V\) is the set of common zeros of \(f_1,...,f_m\) in \(\cc^n\), and it is non-empty, then for any other polynomial \(g \in \mathcal{O}(k^n)\)
    \[g(x) = 0, \forall x \in V \iff g^k \in (f_1,\dots,f_m) \text{ for some } k \in \nn\]
\end{theorem}
We will give a general proof later, but firstly we can see how it makes sense. In above we have shown that \(\sqrt{I} \subset \mathcal{I}(\mathcal{V}(I))\), the other side of inclusion is to ask: when we take a polynomial \(g\) vanishing on \(\mathcal{V}(I)\), what can we say about \(g\)? Clearly, \(I = (f_1,\dots,f_m)\) since we know 
\[\mathcal{V}(\{f_1,\dots,f_m\}) = \mathcal{V}((f_1,\dots,f_m))\]
then the theorem gives the answer when \(k = \cc\), it shows that \(g \in \sqrt{I}\). Hence we have
\[\mathcal{I}(\mathcal{V}(I)) = \sqrt{I} = I ?\]
the first equality holds by the theorem, but the second equality holds under the assumption that \(I\) is a radical ideal.

\begin{colortheorem}[weak Nullstellensatz]
    Let \(K\) be the finitely generated \(k\)-algebra, if \(K\) is a field extension of \(k\), then \(K\) is a finite extension of \(k\).
 \end{colortheorem}

\begin{proof}
    Let \(K\) be generated by \(n\) elements over \(k\), nothing needs to prove, for \(n \geq 1\) we can prove it by induction:\\

    (1) We suppose that for \(n-1\) the statement holds, then for \(n\), there exists \(a_1,\dots,a_n\) in \(K\) such that \(K= k[a_1,...,a_n] = k(a_1)[a_2,...,a_n]\), then by induction \(K\) is a finite extension of \(k(a_1)\) and it suffices to show that \(k(a_1)\) is a finite extension of \(k\), then
    \[[K:k] = [K/k(a_1)] \cdot [k(a_1):k]\]

    (2) \(K\) is finite over \(k(a_1)\), which allows us to find a finite basis \(\{e_1,...,e_m\}\) for \(K\), then for any generator \(a_i\), there exists \(d_{i,k} \in k(a_1)\) such that
    \[a_i  = \sum_{k=1}^{m}d_{i,k}\cdot e_k\]
    and algebra allows us to get \(c_{rs,k}\) such that
    \[e_r \cdot e_s = \sum_{t=1}^m c_{rs,t} \cdot e_t\]
    Hence we can define a finite generated subalgebra \(A = k[d_{i,k},c_{rs,t}]_{i,r,s,k,t} \subset k(a_1)\), then It is not difficult to deduce
    \[K = k[a_1,...,a_n] \subset \oplus_{j=1}^m Ae_j\]
    since each elemet of \(K\) is a \(k(a_1)\)-linear combination of \(e_j\)'s, and each \(a_i\) can be expanded as an \(A\)-linear combination of \(e_j\).

    (2) We finish the proof by absurd: if \(k(a_1)\) is not finite over \(k\), then it is not algebraic over \(k\), i.e. \(k[a_1] \cong k[X]\), then any generator of \(A \subset k(a_1)\) can be written as the form of \(f_i/g_i\), and we can construct a polynoimial coprime with all \(g_i\) by
    \[P = 1 + \prod_{i}g_i\]
    then \(1/P \notin A\), and it shows that \(1/P \notin K\), and then \(K\) is not a field. In fact, if \(1/P \in K\), then there exists \(h_j \in A\) such that \(1/P = \sum_{j=1}^{m}h_j \cdot e_j\), then
    \[1 = \sum_{j=1}^m (P \cdot h_j) \cdot e_j\]
    by direct sum the representation is unique by \(1 = 1+0...+0\), so it shows that \(1/P = h_1 \in A\), which is a contraction.
    
\end{proof}

\begin{remark}
    One point of \(k^n\) determines a maximal ideal of \(k[X_1,...,X_n]\) by
    \[m_a = (X_1-a_1,\dots,X_n-a_n)  \quad \text{if } a = (a_1,\dots,a_n)\]
    clearly that \(k[X_1,\dots,X_n]/m_a \cong k\) is a field, so \(m_a\) is exactly a maximal ideal. Conversely, if \(\mathfrak{m}\) is a maximal ideal of \(k[X_1,...,X_n]\), then \(K = k[X_1,...,X_n]/\mathfrak{m}\) is a finitely generated \(k\)-algebra and it is a field, so \(K \cong k\) if \(k\) is algebraically closed.
\end{remark}

\begin{colorcorollary}
    If \(k\) is algebraically closed, and then

    \begin{enumerate}[(a)]
    \item  For any ideals \(I,J \subset K[X_1,\dots,X_n]\), \(\mathcal{V}(I) = \mathcal{V}(J) \iff \sqrt{I} = \sqrt{J}\).


    \item \textcolor{blue}{(Strong-Nullstellensatz)}  \(\mathcal{I}(\mathcal{V}(I)) = \sqrt{I}\) for any ideals.
    
    \item application \(I \mapsto \mathcal{V}(I)\) and \(V \mapsto \mathcal{I}(V)\) induces a bijection
    \[\{\text{radical ideal of } K[X_1,\dots,X_n]\} \longleftrightarrow \{\text{algebraic variety of } k^n\}\]
    and
    \[\{\text{maximal ideal of } K[X_1,\dots,X_n]\} \longleftrightarrow \{\text{one-point variety of }  k^n\}\]
    \item \(\mathcal{V}(I) = \emptyset \iff I = (1) = k[X_1,\dots,X_n]\).
    \end{enumerate}

\end{colorcorollary}
\begin{proof}
    For the second bijection of (b), it is enough to prove that \(a \mapsto m_a\) is surjective. let \(\mathfrak{m}\) be a maximal ideal, then by remark above we can get a canonical map
    \[\pi : k[X_1,\dots,X_n] \mapsto k[X_1,\dots,X_n]/\mathfrak{m} \cong k\]
    so \(\pi(X_i) = a_i \in k\) for each \(i\), which is enough to prove that \(\mathfrak{m} = m_a\) for \(a = (a_1,\dots,a_n)\).

    \textcolor{blue}{the other results are the consequence of \((a)\), so we prove it firstly:}
    
    (a) \(\Leftarrow\) is clear. Conversely, a radical ideal \(I\) can be written as the intersection of some maximal ideals\textbf{(verify)},
    \[\sqrt{I} = \bigcap_{I \subset \mathfrak{m} } \mathfrak{m}\]
    then by weak nullstellensatz \(\mathcal{V}(\mathfrak{m}) = \{a\}\) for some point of \(k^n\), hence \(I \subset \mathfrak{m}\) if and only if \(a \in \mathcal{V}(I)\), i.e.
    \[\sqrt{I} = \bigcup_{a \in \mathcal{V}(I)} \mathfrak{m}_a\]
    where notice that \(a=(a_1,\dots,a_n) \mapsto (X-a_1,\dots,X-a_n) = \mathfrak{m}_a\) gives a bijection by nullstellensatz. Hence \(V(I) = V(J)\) gives the same radical ideal.

    (b) Set \(S = \mathcal{V}(I)\), then 
    \[\mathcal{V}(\mathcal{I}(S)) = S = \mathcal{V}(I)\]
    the first holds by (e) of proposition \ref{vanish-ideal}, so by part (a) we can get \(\sqrt{\mathcal{I}(S)} = \sqrt{I}\), and then  \(\mathcal{I}(S) = \sqrt{I}\) since \(\mathcal{I}(S)\) is a radical ideal.

    (c) is the immediate consequnce of \((a)\) or \((b)\).



    (d) \(\Leftarrow\) is clear. Conversely we assume that \(I \subsetneq K[X_1,\dots,X_n]\) is a proper ideal, then there exists a maximal ideal \(\mathfrak{m}\) such that \(I \subset \mathfrak{m}\) by Zorn's lemma, then by (a) we can conclude that \(f \in I\) vanishes at some point \(a \in k^n\), hence \(V(I) \neq \emptyset\) which arises a contradiction.
\end{proof}
\begin{remark}
    Strong-Nullstellensatz is a generalization of the original statement of Nullstellensatz (theorem \ref{orginal-Nullstellensatz}).
\end{remark}

\subsection{Zariski topology}

To refine the correspondence about prime ideals, it is better to talk about it in the context of geometry.

\begin{colordefinition}
    The \textbf{Zariski topology} on \(k^n\) is defined by taking the algebraic subsets as all closed subsets, and we write \(\mathbb{A}_k^n\) to stress the topological space \(k^n\) with Zariski topology, and we call it the \textbf{affine n-space} over \(k\).
\end{colordefinition}

such definition is well-defined along the proposition \ref{algberaic-subset}, and a sepcial open subsets can be defined as following:
    \[D(f) := k^n \setminus \mathcal{V}(f) = \{x \in k^n \mid f(x) \neq 0\}\]
where \(f \in k[X_1,...,X_n]\) is a polynomial, and it is called the \textbf{principal open subset} defined by \(f\).

\begin{proposition} \label{principal-open}
    Let \(k\) be a field, then 
    \begin{enumerate}[(a)]
        \item \(D(0) = \emptyset\) and \(D(1) = k^n\).
        \item \(D(f) \cap D(g) = D(fg)\) for any two polynomials \(f,g\).
        \item principal open subsets form a open basis of Zariski topology.
    \end{enumerate}
\end{proposition}

\begin{proof}
    (a) is clear. (b) holds by De-morgan's law
    \[D(f) \cap D(g) = k^n - \mathcal{V}(f) \cup \mathcal{V}(g) = k^n - \mathcal{V}(fg) = D(fg)\]
    (c) Any open subset \(U = k^n - \mathcal{V}(I)\) for some ideal \(I\). \(I = (f_1,\dots,f_m)\) is finitely generated by Noetherian condition of \(k[X_1,...,X_n]\), then each open subset can be written as the union of \textcolor{blue}{finitely many principal open subsets}:
    \[U = k^n - \cap_{i} \mathcal{V}(f_i) = \cup_i D(f_i)\]
\end{proof}

\begin{remark}
    When \(k\) is infinite, \(D(f) = \emptyset\) if and only if \(f = 0\), it is not ture when \(k\) is finite. For example, if \(k = \ff_{p^n}\) and we take \(f(X) = X^{p^n} - X\), it vanishes at all points of \(k\).
\end{remark}

We can investigate the topological properties of Zariski topology, which makes affine space a sepcial topological space.
\begin{itemize}
    \item If \(k\) is finite, then the Zariski topology on \(k^n\) is just the discrete topology.
    \item[] \textcolor{blue}{If \(k\) is infinite.}

    \item  When \(n = 1\), the Zariski topology on \(k\) is just the cofinite topology: for any finite subset \(F = \{c_1,...,c_n\}\subset k\), we can find that \(f(X) = (X-c_1)\cdots(X-c_n)\) vanishes exactly on \(F\), hence \(V(f) = F\), which shows that any finite subset is closed, and there does not exist a polynomial vanishing at inifinitely many points in \(k[X]\), hence all possible closed sets are
    \[\mathcal{F} = \{k,\emptyset,\text{finite subsets}\}\]
    which is exactly the cofinite topology, it is a natural example to understand the Zariski topology. \textbf{Pay attention} that when \(n\geq 2\), the Zariski topology is not the cofinite topology, a example is that 
    \[V(X_1) = \{(0,a_2,\dots,a_n) \mid a_i \in k\}\]\\
    it is an infinite closed subset of \(k^n\), hence 
    \[\text{cofinite topology} < \text{Zariski topology} < \text{Euclidean topology}\]
    when “<" means “is coarser than".

    \item Hilbert's basis theorem shows that \(k[X_1,...,X_n]\) is Noetherian, hence the Zariski topology on \(k^n\) is \textcolor{blue}{a Noetherian topological space}: any ascending chain of open subsets \(U_1 \subset U_2 \subset \cdots\) will stop at \(U_N\) for some integer \(N > 0\).
    
    For any chain of algebraic subsets \(F_1 \supseteq F_2 \supseteq \cdots\), the vanishing ideals allow us to get an increasing chain of ideals
    \[\mathcal{I}(F_1) \subseteq \mathcal{I}(F_2) \subseteq \cdots\]
    so by noetherian condition we can finish the proof.

    \item \textcolor{blue}{The affine space is a \(T_1\)-space (any single point is closed)}: Indeed, each point \(a \in k^n\) can be seen as a algebraic subset by
    \[\{a\} = \mathcal{V}(X_1 - a_1,...,X_n - a_n)\]
    Howerver, \textbf{the Zariski topology is not \(T_2\) (Hausdorff)}, consider any two points \(a,b \in k^n\), and two non-empty set principal open subsets \(D(f), D(g)\) containing them respectively, then
    \[D(f) \cap D(g) = \emptyset \iff D(fg) = \emptyset \iff fg = 0\]
    \(k[X_1,...,X_n]\) is a domain, so \(f = 0\) or \(g = 0\), which leads to  a contradiction.
    
    \item \textcolor{blue}{The Zariski topology is quasi-compact}: any open cover has a finite subcover (see the proof of  Proposition \ref{principal-open} (c)).
    
    \item \textcolor{blue}{The Zariski topology is connected}: Suppose that \(k^n = U \cup V\) are two non-empty disjoint open subsets, then the open basis allows to choose non-empty open sets \(D(f)\) and \(D(g)\) such that \(D(f)\cap D(g)\), then disjoint condition shows that \(D(fg) = \emptyset\), hence \(fg\) is nilpotent, then \(D(f) =\emptyset\) or \(D(g) = \emptyset\), which leads to contradiction.
\end{itemize}

Another condition related to the prime ideals is the irreducibility of algebraic subsets, which is defined as following:
\begin{colordefinition}
    A topological space \(X\) is called \textbf{irreducible} if it satisfies the following equivalent conditions:
    \begin{enumerate}[(a)]
        \item \(X\) is not the union of two proper closed subsets.
        \item Any two non-empty open subsets of \(X\) have non-empty intersection.
        \item Any non-empty open subset of \(X\) is dense in \(X\).
    \end{enumerate}
\end{colordefinition}

\begin{proof}
    
\end{proof}

\begin{remark}
    It is a stronger condition than connectedness, for example, the union of two lines \(X = 0\) and \(Y = 0\) in \(\rr^2\), it can be seen as the algebraic subset \(\mathcal{V}(XY) \) which is not irreducible since it is the union of two algebraic subsets \(\mathcal{V}(X)\) and \(\mathcal{V}(Y)\).

    It is connected under the Zariski topology, to see that we notice \(\mathcal{V}(X) \cap \mathcal{V}(Y) = \{(0,0)\}\) is not empty, while \(\mathcal{V}(X)\) and \(\mathcal{V}(Y)\) homeomorphic to \(\mathbb{A}\) which is connected, so the union is connected.
\end{remark}

\begin{colorproposition}
    Let \(k\) be an algebraically closed field, then there exists a correspondence by \(I \mapsto \mathcal{V}(I)\) and \(V \mapsto \mathcal{I}(V)\):
    \[\{\text{prime ideal of } k[X_1,\dots,X_n]\} \longleftrightarrow \{\text{irreducible algebraic variety of } k^n\}\]
    such a varity is called \textbf{affine variety}.
\end{colorproposition}
\begin{proof}
    Let \(I\) be a prime ideal of \(k[X_1,\dots,X_n]\), then we need to show that \(V(I)\) is irreducible. If \(V(J)\) and \(V(L)\) are two algebraic subsets such that
    \[V(I) = V(J) \cup V(L) = V(JL)\]
    then by the correspondence of radical ideals, we have \(\sqrt{I} = \sqrt{JL}\), \(I\) is prime so it is radical, hence \(I = JL\), then by the definition of prime ideal, \(J \subset I\) or \(L \subset I\), hence \(V(I) \subset V(J)\) or \(V(I) \subset V(L)\), which shows that \(V(I)\) is irreducible.\\

    Let \(V\) be an irreducible algebraic subset of \(k^n\), then we need to show that \(I(V)\) is a prime ideal. For any \(f,g \in K[X_1,\dots,X_n]\) such that
    \[fg \in I(V) \iff V \subset V(fg) = V(f) \cup V(g)\]
    then by irreducibility of \(V\), we have \(V \subset V(f)\) or \(V \subset V(g)\), hence \(f \in I(V)\) or \(g \in I(V)\), which shows that \(I(V)\) is prime.
\end{proof}

Noetherian space allows us to define the dimension of topological space by the length of chain, which can be formulated as following:
\begin{definition}
    The dimension of a topological space \(X\) is defined by the largest integer \(n\) such that there exists a chain of affine varieties of \(X\):
    \[Z_0 \subsetneq  Z_1 \subsetneq \cdots \subsetneq Z_n \]

    A quasi-affine variety is an open subset of an affine variety, and the dimension of a (quasi-)affine variety is defined to be its dimension as the subspace of affine space.
\end{definition}



\subsection{Coordinate ring and functor \(\mathcal{O}\)}

This section is to give a complete description about the category of affine varieties over \(k\).

\begin{definition}
    Let \(V \subset k^n\) be an algebraic subset, then a function \(f: V \to k\) is called a \textbf{polynomial function} on \(V\) if there exists a polynomial function \(\tilde{f} \in \mathcal{O}(k^n)\) such that \(f = \tilde{f}|_V\). We denote the set of all polynomial functions on \(V\) by \(\mathcal{O}(V)\).\\
\end{definition} 

The restriction induces a natural surjective morphism of \(k\)-algebras
\[r: \mathcal{O}(k^n) \to \mathcal{O}(V), \quad f \mapsto f|_V\]
with the kernel \(\ker r = \mathcal{I}(V)\), hence it gives an isomorphism \(\oo(V) \cong \oo(k^n)/\mathcal{I}(V)\)
and we call \(\mathcal{O}(V)\) the \textbf{coordinate ring} of the variety \(V\), the following corollary reflects the name of "coordinate", it reflects all the information of the variety:

\begin{proposition}[UPQ-Coordinate ring] $ \\$
    Let \(V \subset k^n\) be an algebraic subset \textcolor{blue}{with \(k\) algebraically closed}, then  there exists a natural bijection 
    \[V \longleftrightarrow \Hom_{k-alg}(\mathcal{O}(V),k)\]
    by mapping a point \(a \in V\) to the evaluation map \(\mathrm{ev}_a\):
    \[\mathrm{ev}_a: \mathcal{O}(V) \to k, \quad f \mapsto f(a)\]
\end{proposition}
\begin{proof}

\end{proof}

\begin{example}
    We should notice that the coordinate ring actually reflects the geometric information of variety in the way of algebra, for example we condiser two algebraic subsets
\[V_1 = \{(x,0) \in k^2 \mid x \in k\} \quad V_2 = \{(x,x,x) \in k^3 \mid x\in k\}\]
clearly that they are all strignt lines in different affine spaces, and we can calculate their coordiate rings:
\[\mathcal{O}(V_1) \cong k[X,Y]/(Y) \cong k[X], \quad \mathcal{O}(V_2) \cong k[X,Y,Z]/(X-Y,Y-Z) \cong k[X]\]
in some sense, \(V_1\) and \(V_2\) are all curves, which can be reflected by the dimension of coordinate rings.
\end{example}

\begin{proposition}
    Let \(k\) be an \textcolor{blue}{algebraically closed field}, and \(V \subset \mathbb{A}_k^n\) be an affine variety, then dimension of \(V\) is equal to the Krull dimension of \(\mathcal{O}(V)\).
\end{proposition}

\begin{proof}
    
\end{proof}

\begin{definition}
    Let \(V \subset k^n\) and \(W \subset k^m\) be two algebraic subsets, then a map \(\varphi: V \to W\) is called a \textbf{polynomials} if it is a restriction of a polynomials map \(k^n \to k^m\)
    \[(x_1,\dots,x_n) \mapsto (f_1(x_1,\dots,x_n),\dots,f_m(x_1,\dots,x_n))\]
    where \(f_i \in \mathcal{O}(k^n)\).
\end{definition}

Indeed such a map is continous under Zariski topology \textbf{(verify!)}, which allows us to construct a sub-category of \(\mathbf{Top}\), called t\textbf{he category of affine algebraic varieties over \(k\)}:
\[\cat{\mathbf{Aff}_k}{\text{algebraic variety } V \subset k^n}{\text{polynomial functions } f: V\subset k^n \to W \subset k^m}\]

with this definition, then as a functor \(\mathcal{O}\) maps each object to a coordinate ring, exactly a \(k\)-algebra of type fini, here we denote it by \(\mathbf{Alg}_k^{\text{fini}}\).\\

On the level of morphisms, each polynoimial function \(\varphi:k^n \to k^m\) induces a morphism of \(k\)-algebra like dual map:
\[\varphi^*: \mathcal{O}(k^m) \to \mathcal{O}(k^n), \quad f \mapsto f\circ \varphi \]
hence furthermore let \(\varphi|V:V \to W\) be the restriction, then it induces a morphism of coordiante rings by quoitent:
% https://q.uiver.app/#q=WzAsNCxbMCwwLCJcXG1hdGhjYWx7T30oa15tKSJdLFsyLDAsIlxcbWF0aGNhbHtPfShrXm4pIl0sWzAsMiwiXFxtYXRoY2Fse099KFcpIl0sWzIsMiwiXFxtYXRoY2Fse099KFYpIl0sWzAsMSwiXFx2YXJwaGleKiJdLFswLDIsIlxccGlfVyIsMl0sWzEsMywiXFxwaV9WIl0sWzIsMywiIiwyLHsic3R5bGUiOnsiYm9keSI6eyJuYW1lIjoiZGFzaGVkIn19fV1d
\[\begin{tikzcd}
	{\mathcal{O}(k^m)} && {\mathcal{O}(k^n)} \\
	\\
	{\mathcal{O}(W)} && {\mathcal{O}(V)}
	\arrow["{\varphi^*}", from=1-1, to=1-3]
	\arrow["{\pi_W}"', from=1-1, to=3-1]
	\arrow["{\pi_V}", from=1-3, to=3-3]
	\arrow[dashed, from=3-1, to=3-3]
\end{tikzcd}\]
It is easy to check:
\begin{align*}
    \varphi^*(f) \in I(V) &\iff f \circ \varphi \in I(V) \\
    &\iff f(\varphi(a)) = 0, \forall a \in V\\
    &\iff f(b) = 0, \forall b \in W \quad\\
    &\iff f \in I(W)
\end{align*}
hence it induces a unique morphism of \(k\)-algebras
\[\varphi^*: \mathcal{O}(W) \to \mathcal{O}(V), \quad f+I(W) \mapsto f \circ \varphi + I(V)\]

Hence we finish the consturction of functor completely
\[\mathcal{O}: \mathbf{Aff}_k \to \mathbf{Alg}_k^{\text{fini}}\]
without the assumption that \(k\) is algebraically closed, the functor will not be an equivalence, actually we need another condition to ensure that, which makes us can give a one-side correpsondence completely:
\begin{theorem}
    Let \(k\) be an algebraically closed field, then the functor \(\mathcal{O}\) gives a contravariant equivalence to the reduced \(k\)-algebras of type fini:
    \[\mathbf{Aff}_k \cong (\mathbf{Alg}_k^{\text{fini, red}})^{\text{op}}\]
\end{theorem}

